\begin{definicija}
  Naj bo $F: I \times \R \times \R \to \R$ zvezna funkcija in $I$ interval v
  $\R$.
  \pojem{Navadna diferencialna enačba prvega reda} je enačba oblike $F(x, y(x),
  y'(x)) = 0$, kjer je $y(x)$ neka funkcija.
  Rešitev enačbe je vsaka funkcija $y_r(x) : I \to \R$, za katero velja enačba.
\end{definicija}

\begin{opomba}
  NDE $n$-tega reda definiramo podobno kot enačbo oblike
  \[
    F(x, y, y', y'', \ldots, y^{(n)}) = 0.
  \]
\end{opomba}

\begin{opomba}
  Smiselno je opazovati tudi enačbe, kjer je $F = (F_1, \ldots, F_m)$ vektorska
  funkcija.
  Temu pravimo \pojem{sistem NDE}.
\end{opomba}

\begin{opomba}
  Naj bo $y^{(n)} = F(x, y, y', \ldots, y^{(n-1)})$ enačba reda $n$.
  Ta enačba je ekvivalentna primernemu sistemu $n \times n$ prvega reda;
  definirajmo $y_1 = y$, $y_2 = y'$, \ldots, $y_{n-1} = y^{n-2}$.
  Tedaj dobimo enačbo $y_n' = F(x, y_1, \ldots, y_n)$.
\end{opomba}

\begin{definicija}
  \pojem{Integralska krivulja} $\gamma$ vektorskega polja $F: \Omega \subseteq
  \R^n \to \R^n$ skozi točko $x_0 \in \Omega$ je krivulja $\gamma : [0, b) \to
  \Omega$, za katero velja
  \begin{itemize}
  \item v vsaki točki $t$ je $\dot{\gamma}(t) = F(\gamma(t))$,
  \item $\gamma(0) = x_0$.
  \end{itemize}
\end{definicija}

\vprasanje{Definiraj integralske krivulje.}

Če prvi pogoj iz definicije zapišemo v koordinatah,
\[
  \begin{bmatrix}
    \dot{x}_1 \\ \vdots \\ \dot{x}_n
  \end{bmatrix} (t)
  =
  \begin{bmatrix}
    F_1(x_1, \ldots, x_n) \\ \vdots \\ F_n(x_1, \ldots, x_n)
  \end{bmatrix},
\]
dobimo sistem $n$ NDE prvega reda z $n$ neznankami.
Ta sistem ni eksplicitno odvisen od $t$; takim sistemom pravimo \pojem{avtonomni
sistemi}.
Pokazali bomo, da za vsako izbiro $x_0$ obstajata interval $[0, a)$ in krivulja
$\gamma$, za katero veljata pogoja v definiciji.

Vsak neavtonomen sistem lahko prepišemo v avtonomnega, z uvedbo nove odvisne
spremenljivke $v(t) = t$.
Dobimo nov sistem
\begin{gather*}
  \dot{v} = 1, \\
  \dot{x}_1 = F_1(v, x_1, \ldots, x_n), \\
  \vdots \\
  \dot{x}_n = F_n(v, x_1, \ldots, x_n).
\end{gather*}
Partikularna rešitev tega sistema je tedaj integralska krivulja vektorskega
polja $\vec{F}(v, \vec{x})$ v \pojem{razširjenem faznem prostoru} $\R \times
\Omega$ ($\Omega$ je običajen fazni prostor), ki ustreza primernemu začetnemu
pogoju.

\vprasanje{Kako spremenimo neavtonomni sistem v avtonomnega?}

